
%%%%%%%%%%%%%%%%%%%%%%%%%%%%%%%%%%%%%%%%%
% Beamer Presentation
% LaTeX Template
% Version 1.0 (10/11/12)
%
% This template has been downloaded from:
% http://www.LaTeXTemplates.com
%
% License:
% CC BY-NC-SA 3.0 (http://creativecommons.org/licenses/by-nc-sa/3.0/)
%
%%%%%%%%%%%%%%%%%%%%%%%%%%%%%%%%%%%%%%%%%

% ----------------------------------------------------------------------------------------
% PACKAGES AND THEMES
% ----------------------------------------------------------------------------------------

\documentclass{beamer}

\usepackage{tikz}
\usetikzlibrary{trees}
\usetikzlibrary{cd}

\usepackage{graphicx} % Allows including images
\usepackage{mathpartir} % Allows including images
\graphicspath{{media/}}
% \usepackage{subfig} % Allows including images
% \usepackage{booktabs} % Allows the use of \toprule, \midrule and \bottomrule in tables
\usepackage[]{bm}
\usepackage[backend=biber,
style=authoryear-comp,
]{biblatex}
\addbibresource{../sources.bib}

\mode<presentation> {

  % The Beamer class comes with a number of default slide themes
  % which change the colors and layouts of slides. Below this is a list
  % of all the themes, uncomment each in turn to see what they look like.

  % \usetheme{default}
  % \usetheme{AnnArbor}
  % \usetheme{Antibes}
  % \usetheme{Bergen}
  % \usetheme{Berkeley}
  % \usetheme{Berlin}
  % \usetheme{Boadilla}
  % \usetheme{CambridgeUS}
  % \usetheme{Copenhagen}
  % \usetheme{Darmstadt}
  % \usetheme{Dresden}
  % \usetheme{Frankfurt}
  % \usetheme{Goettingen}
  % \usetheme{Hannover}
  % \usetheme{Ilmenau}
  % \usetheme{JuanLesPins}
  % \usetheme{Luebeck}
  % \usetheme{Madrid}
  % \usetheme{Malmoe}
  % \usetheme{Marburg}
  % \usetheme{Montpellier}
  % \usetheme{PaloAlto}
  % \usetheme{Pittsburgh}
  % \usetheme{Rochester}
  % \usetheme{Singapore}
  % \usetheme{Szeged}
  \usetheme{Warsaw}

  % As well as themes, the Beamer class has a number of color themes
  % for any slide theme. Uncomment each of these in turn to see how it
  % changes the colors of your current slide theme.

  % \usecolortheme{albatross}
  % \usecolortheme{beaver}
  % \usecolortheme{beetle}
  % \usecolortheme{crane}
  % \usecolortheme{dolphin}
  % \usecolortheme{dove}
  % \usecolortheme{fly}
  % \usecolortheme{lily}
  % \usecolortheme{orchid}
  % \usecolortheme{rose}
  % \usecolortheme{seagull}
  % \usecolortheme{seahorse}
  % \usecolortheme{whale}
  % \usecolortheme{wolverine}
  % \usecolortheme{spruce}
  % \usecolortheme{structure}
  % \setbeamertemplate{footline} % To remove the footer line in all slides uncomment this line
  \setbeamertemplate{footline}[frame number] % To replace the footer line in all slides with a simple slide count uncomment this line

  \setbeamertemplate{navigation symbols}{} % To remove the navigation symbols from the bottom of all slides uncomment this line
}

\newcommand{\email}[1]{\href{mailto:#1}{\texttt{#1}}}
\renewcommand{\note}[1]{{\color{orange}#1}}
\newcommand{\faded}[2][35]{\textcolor{fg!#1}{#2}}

\DeclareMathOperator{\ifs}{\mathtt{if}}
\DeclareMathOperator{\whiles}{\mathtt{while}}
\DeclareMathOperator{\skips}{\mathtt{skip}}

\DeclareMathOperator{\fold}{\mathtt{fold}}
\DeclareMathOperator{\update}{\mathtt{update}}
\DeclareMathOperator{\true}{\mathtt{true}}
\DeclareMathOperator{\false}{\mathtt{false}}

% syntax for traces
\newcommand{\pc}[1]{\ensuremath{pc(#1)}}
\newcommand{\accC}[2][V_0]{\ensuremath{#2\Downarrow_{#1}}}
\newcommand{\accS}[2][]{\ensuremath{#2\downarrow_{#1}}}

% ----------------------------------------------------------------------------------------
% TITLE PAGE
% ----------------------------------------------------------------------------------------

\title[Concurrent SE in Coq]{Concurrent Symbolic Execution with Trace Semantics in Coq}
\author[Å. A. A. Kløvstad]{Åsmund Aqissiaq Arild Kløvstad \\he/him\\ \email{aaklovst@ifi.uio.no}}
\institute[UiO]{\includegraphics[height=2cm]{03_uio_full_logo_eng_pos} \\
Department of Informatics}
\date[PLACES 2023]{PLACES 2023}
% {Programming Language Approaches to Concurrency- and Communication-cEntric Software 2023} % Date, can be changed to a custom date

\begin{document}

\begin{frame}
  \titlepage%
\end{frame}

% \begin{frame}
%   \tableofcontents%
% \end{frame}
%
\begin{frame}{Motivation}
  Symbolic execution for concurrent programs
  \begin{itemize}[<+->]
    \item Scaling problems (and solutions)
    \item Formal model -- correctness and completeness
  \end{itemize}
  \vfill{}
  \uncover<+->{This is work in progress!}
\end{frame}

\section{Symbolic Execution Formally Explained}

\subsection{Symbolic Execution}

\begin{frame}{What (and why) is Symbolic Execution}
  \begin{columns}
    \begin{column}{.4\textwidth}
      \begin{itemize}
        \item<+-> Symbolic values
        \item<+-> Path condition
        \item<+-> Uses:
              \begin{itemize}
                \item Proving (partial) correctness
                \item Static analysis
                \item (Directed) testing
              \end{itemize}
        \item<+-> Challenges:
              \begin{itemize}
                \item Formalism
                \item \alert{Scaling with concurrency}
                \item \faded{Non-termination}
              \end{itemize}
      \end{itemize}
    \end{column}
    \begin{column}<2->{.6\textwidth}
      \begin{block}{Ex. \texttt{while x > y \{x := x - y\}}}
        \begin{figure}
          \begin{tikzpicture}[grow via three points={%
              one child at (0,-1.5) and two children at (-1.5,-1.25) and (1.5,-1.25)}]
            \node at (0,0) [align=center]
            {$\mathtt{x} \mapsto x_{0}, \mathtt{y} \mapsto y_{0}$}
            child {node [align=center]
              {$\mathtt{x} \mapsto x_{0}, \mathtt{y} \mapsto y_{0}$ \\ $x_{0} > y_{0}$}
              child {node [align=center]
                {$\mathtt{x} \mapsto x_{0} - y_{0}, \mathtt{y} \mapsto y_{0}$ \\ $x_0 > y_0$}
                child {node [align=center] {\ldots}}
                child {node [align=center]
                  {$\mathtt{x} \mapsto x_{0} - y_{0}, \mathtt{y} \mapsto y_{0}$ \\ $x_{0} > y_{0} \land x_{0} - y_{0} \leq y_{0}$}}
              }}
            child {node [align=center] {$\mathtt{x} \mapsto x_{0}, \mathtt{y} \mapsto y_{0}$ \\ $x_{0} \leq y_{0}$}};
          \end{tikzpicture}
        \end{figure}
      \end{block}
    \end{column}
  \end{columns}
\end{frame}

\subsection{Formal Framework}

\begin{frame}{SE Formally Explained~\parencite{boer2021}}
  \begin{columns}[t]
    \begin{column}{.6\textwidth}
      \begin{itemize}[<+->]
        \item Expressions
              \begin{itemize}
                \item<.-> With standard booleans ops
              \end{itemize}
        \item A small imperative language
        \item Substitutions $\sigma$: Var $\rightarrow$ Expr
              \begin{itemize}
                \item<.-> Application $e\sigma$: Expr
                \item<.-> Identity $id : x \mapsto x$
              \end{itemize}
        \item Path condition $\phi$: boolean expression
        \item Valuations $V$: Var $\rightarrow$ Val
              \begin{itemize}
                \item<.-> Evaluation $V(e)$: Val
              \end{itemize}
        \item $V \circ \sigma$: Valuation
              \begin{itemize}
                \item<.-> Abstraction/concretization
                \item<.-> $V \circ id = V$
              \end{itemize}
      \end{itemize}
    \end{column}

    \begin{column}{.4\textwidth}
      \begin{block}<1->{Expressions}
        \vspace{-3ex}
        \begin{align*}
          e ::&= x \\
              & \mid op(e_{1},\ldots, e_{n})
        \end{align*}
      \end{block}
      \begin{block}<2->{Statements}
        \vspace{-3ex}
        \begin{align*}
          s ::&= x := e\\
              & \mid s~;~s\\
              & \mid \ifs~e~\{s\}~\{s\}\\
              & \mid \whiles~e~\{s\}
        \end{align*}
      \end{block}
      \begin{block}{V-Abstraction}<6->
        $\sigma$ V-abstracts $V'$ if $V \circ \sigma = V'$
      \end{block}
    \end{column}
  \end{columns}
\end{frame}

\newcommand{\Sstep}[2]{\ensuremath{#1 \rightarrow #2}}
\newcommand{\Sstar}[2]{\ensuremath{#1 \rightarrow^{*} #2}}
\newcommand{\Cstep}[3][]{\ensuremath{#2 \bm{\Rightarrow}_{#1} #3}}
\newcommand{\Cstar}[3][]{\ensuremath{#2 \bm{\Rightarrow^{*}}_{#1} #3}}

\begin{frame}{SE Formally Explained: Semantics}
  \begin{block}<+->{Symbolic $\Sstar{(s, \sigma, \phi)}{(s', \sigma', \phi')}$}
    \vspace{-3ex}
    \begin{align*}
      &\Sstep{(x := e~;~s, \sigma, \phi)}{(s, \sigma[x := e\sigma], \phi)}\\
      &\Sstep{(\whiles~e~\{s\}~;~s', \sigma, \phi)}{(s~;\whiles~e~\{s\}~;~s', \sigma, \phi\land e\sigma)}\\
      &\Sstep{(\whiles~e~\{s\}~;~s', \sigma, \phi)}{(s', \sigma, \phi\land \neg e\sigma)}\\
      &\ldots
    \end{align*}
  \end{block}
  \begin{block}<+->{Concrete $\Cstar{(s, V)}{(s', V')}$}
    \vspace{-3ex}
    \begin{align*}
      &\Cstep{(x := e~;~s, V)}{(s, V[x := V(e)])}\\
      &\Cstep{(\whiles~e~\{s\}~;~s', V)}{(s', V)}, \; \text{if V(e) = false}\\
      &\ldots
    \end{align*}
  \end{block}
\end{frame}

\begin{frame}{SE Formally Explained: Correctness \& Completeness}
  \uncover<1->{
    \begin{theorem}[Correctness]
      If $\Sstar{(s, id, true)}{(s', \sigma, \phi)}$ and $V(\phi)$ = true, then $\Cstar{(s, V)}{(s, V \circ \sigma)}$
    \end{theorem}
    Symbolic executions with satisfiable path conditions have corresponding concrete executions
  }
  \vfill
  \uncover<2->{
    \begin{theorem}[Completeness]
      If $\Cstar{(s, V_{0})}{(s', V)}$ there exists $\Sstar{(s, id, true)}{(s', \sigma, \phi)}$ such that $V = V_{0} \circ \sigma$ and $V(\phi)$ = true
    \end{theorem}
    Concrete executions have symbolic counterparts with satisfied path condition
  }
\end{frame}

\subsection{Coq Mechanization}

\begin{frame}{SE Formally Formalized}
  Implementing de Boer and Bonsangue's work in Coq
  \begin{columns}
    \begin{column}{.5\textwidth}
      \begin{itemize}[<+->]
        \item Generic context machinery
              \begin{itemize}
                \item \texttt{C : Stmt -> Stmt}
                \item<.-> \texttt{is\_context}
                \item Head reduction relation
              \end{itemize}
        \item Stepwise transitive closure
        \item Proofs by induction + case analysis
        % \item Typed expressions
        % \item The id-substitution
      \end{itemize}
    \end{column}
    \begin{column}{.5\textwidth}
      \begin{block}{Contexts}<2->
        \vspace{-3ex}
        \begin{align*}
          C ::&= \square \mid C~;~s\\
          s ::&= \faded{\ldots}\\
              &\mid \skips
              % an example, but it's ugly
              % (\square;x:=42)[y:=7] &\triangleq y:=7;x:=42
        \end{align*}
      \end{block}
      \begin{block}{Head Reduction}<3->
        \vspace{-3ex}
        \begin{align*}
          \infer{\Sstep{(s, X)}{(s', X')}}{\Sstep{(C~s, X)}{(C~s', X')}}
        \end{align*}
      \end{block}
    \end{column}
  \end{columns}
  \vfill{}
  \uncover<+->{Ask me about Coq after!}
\end{frame}

\section{Extending the Framework}
\subsection{Parallelism}

\begin{frame}{Parallelism}
  \begin{columns}[t]
    \begin{column}{.6\textwidth}
      \begin{itemize}[<+->]
        \item Parallel composition
        \item Shared memory
        \item Extra context
        \item Correctness and completeness still hold!
      \end{itemize}
      \begin{block}<3->{Contexts}
        \vspace{-3ex}
        \begin{align*}
          C ::&= \faded{\square \mid C~;~s}\\
              &\faded[100]{\mid C~|\!|~s}\\
              &\faded[100]{\mid s~|\!|~C}
        \end{align*}
      \end{block}
    \end{column}
    \begin{column}{.4\textwidth}
      \begin{block}<1->{Statements}
        \vspace{-3ex}
        \begin{align*}
          s ::&= \faded{x := e}\\
              &\faded{\mid s~;~s}\\
              &\faded[100]{\mid s~|\!|~s}\\
              &\faded{\mid \ifs~e~\{s\}~\{s\}}\\
              &\faded{\mid \whiles~e~\{s\}}\\
              &\faded{\mid \skips}
        \end{align*}
      \end{block}
    \end{column}
  \end{columns}
\end{frame}

\subsection{Trace Semantics}

\begin{frame}{Why Trace Semantics}
  \begin{itemize}[<+->]
    \item Expressive -- consider \texttt{x := x + 1 |\!| x := x + 2}
    % \[\sigma[x := x + 1][x := x + 2] = \sigma[x := x + 3] = \sigma[x := x + 2][x := x + 1]\]
    % but $(x := x + 1) \curvearrowright (x := x + 2) \ne (x := x + 2) \curvearrowright (x := x + 1)$
    \item Compositional
    \item Partial Order Reduction~(SymPaths,~\cite{boer2020sympaths})
  \end{itemize}
\end{frame}

\begin{frame}{Traces}
  \begin{columns}
    \begin{column}{.5\textwidth}
      \begin{itemize}[<+->]
        \item Symbolic events: assignments and boolean guards
        \item Concrete events: assignments
        \item Final states
        \item Path condition
        \item Abstraction
      \end{itemize}
    \end{column}
    \begin{column}{.5\textwidth}
      \begin{block}{Traces}
        \vspace{-3ex}
        \begin{align*}
          \uncover<1->{\tau^{S} ::&= [~] \mid \tau^{S} :: (x := e) \mid \tau^{S} :: e\\}
          \uncover<2->{\tau^{C} ::&= [~] \mid \tau^{C} :: (x := e)}
        \end{align*}
      \end{block}
      \begin{block}<3->{}
        \vspace{-3ex}
        \begin{align*}
          \uncover<3->{\accC[V]{\tau^{C}} &\triangleq \fold~\update~V~\tau^{C}\\
          \accS[\sigma]{\tau^{S}} &\triangleq \fold~\update~\sigma~\tau^{S}\\}
          \uncover<4->{\pc{\tau^{S}} &\triangleq \fold~\land~\true~\tau^{S} \\
                                          &\text{(sort of)}}
        \end{align*}
      \end{block}
      \begin{block}<5->{Trace Abstraction}
        $\tau^{S}$ V-abstracts $\tau^{C}$ if $\accC[V]{\tau^{C}} = V \circ \accS[id]{\tau^{S}}$
      \end{block}
    \end{column}
  \end{columns}
\end{frame}

\begin{frame}{Transition System}
  \begin{block}{Symbolic $\Sstar{(s, \tau^{S})}{(s', \tau'^{S})}$}<+->
    \vspace{-3ex}
    \begin{align*}
      &\Sstep{(x := e, \tau^{S})}{(\skips, \tau^{S} :: (x :=  e))}\\
      &\Sstep{(\ifs~e~\{s_{1}\}\{s_{2}\}, \tau^{S})}{(s_{1}, \tau^{S} :: e)}\\
      &\Sstep{(\ifs~e~\{s_{1}\}\{s_{2}\}, \tau^{S})}{(s_{2}, \tau^{S} :: \neg e)}\\
      &\ldots
    \end{align*}
  \end{block}
  \begin{block}{Concrete $\Cstar[V]{(s, \tau^{C})}{(s', \tau'^{C})}$}<+->
    \vspace{-3ex}
    \begin{align*}
      &\Cstep[V]{(x := e, \tau^{C})}{(\skips, \tau^{C} :: (x := e))}\\
      &\Cstep[V]{(\ifs~e~\{s_{1}\}\{s_{2}\}, \tau^{C})}{(s_{1}, \tau^{C})}, \text{if } \accC[V]{\tau^C} (e) = \true\\
      &\ldots
    \end{align*}
  \end{block}
\end{frame}

\begin{frame}{Bisimulation}
  \begin{block}{Bisimulation}
    \begin{figure}
      \begin{tikzcd}[ampersand replacement=\&]% <--- otherwise Beamer breaks without [fragile]
        {(s, \tau_{0}^{S})} \arrow[rr] \arrow[dd, "V-abstracts"', no head, dashed] \&  \& {(s', \tau^{S})} \arrow[dd, "V-abstracts", no head, dashed] \\
        \&  \&                                                             \\
        {(s, \tau_{0}^{C})} \arrow[rr, Rightarrow]                                \&  \& {(s', \tau^{C})}
      \end{tikzcd}
    \end{figure}
  \end{block}
  \begin{itemize}[<+->]
    \item If $\tau^{S}$ has satisfiable path condition, then there is a concrete step
    \item If there is a concrete step, then $\tau^{S}$ has satisfiable path condition
    \item Correctness and completeness follow
  \end{itemize}
\end{frame}

\section{Summary}

% \begin{frame}{Future Work}
%   \begin{itemize}[<+->]
%     \item \alert{Partial Order Reduction}
%     \item Composition
%     \item Generalizations
%           \begin{itemize}
%             \item Language features
%             \item Semantics / trace events
%             \item Different notions of abstraction
%           \end{itemize}
%   \end{itemize}
% \end{frame}

\begin{frame}{Summary}
  \begin{itemize}[<+->]
    \item Symbolic Execution Formally Explained, Formalized
    \item Extended with:
          \begin{itemize}
            \item<.-> Parallel Composition
            \item<.-> Trace Semantics
            \item<.-> Contexts
          \end{itemize}
    \item Future directions
          \begin{itemize}
            \item<.-> Partial Order Reduction
            \item<.-> Compositionality
            \item<.-> Language features
          \end{itemize}
  \end{itemize}
  \vfill{}
  \uncover<+->{\center{Thank you, questions?}}
\end{frame}

\appendix
\begin{frame}{References}
  \printbibliography{}
\end{frame}

\begin{frame}{Bisimulation}
  \begin{theorem}[Bisimulation]
    If $\accC[V]{\tau^{C}_{0}} = V \circ\accS{\tau^{S}_{0}}$ and $V(\pc{\tau^{S}_{0}}) = \true$ then:
    \begin{itemize}
      \item if $\Sstep{(s, \tau^{S}_{0})}{(s', \tau^{S})}$ and $V(\pc{\tau^{S}}) = \true$, there exists $\Cstep[V]{(s, \tau^{C}_{0})}{(s', \tau^{C})}$
      \item if $\Cstep[V]{(s, \tau^{C}_{0})}{(s', \tau^{C})}$, there exists $\Sstep{(s, \tau^{S}_{0})}{(s', \tau^{S})}$
            with true path condition
    \end{itemize}
  \end{theorem}
\end{frame}

\begin{frame}{Correctness \& Completeness Follow}
  \begin{corollary}[Correctness]
    If $\Sstar{(s, \tau^{S}_{0})}{(s', \tau)}$, $\tau^{C}_{0}$ V-abstracts $\tau^{S}_{0}$, and $V(\pc{\tau})$ = true, there exists $\Cstar[V]{(s, \tau^{C}_{0})}{(s', \tau^{C})}$ such that $\tau^{S}$ V-abstracts $\tau^{C}$
  \end{corollary}
  \vfill
  \begin{corollary}[Completeness]
    If $\Cstar[V]{(s, \tau_{0}^{C})}{(s', \tau_{0}^{C})}$ there exists $\Sstar{(s, \tau_{0}^{S})}{(s', \tau^{S})}$ such that:
    \begin{itemize}
      \item $V(\pc{\tau^S})$ = true
      \item $\tau_{0}^S$~V-abstracts~$\tau_{0}^C$
      \item $\tau^S$~V-abstracts~$\tau^C$
    \end{itemize}
  \end{corollary}
\end{frame}

\end{document}
